% GENERATED FILE. Edit canonical lessons, then run npm run generate:books.
% canonical-source-hash: fnv1a64:4af804c823f0c067
% canonical-lessons: JA-C65-koi, JA-C65-semi, JA-C65-tonbo, JA-C65-imo, JA-C65-gobou

\chapter{Carp, Cicada, Dragonfly, Potato, Burdock root}
\label{ch:ja-carp-cicada-dragonfly-potato-burdock-root}
\hlchaptermodality{65}

\begin{chapteropening}
\textbf{By the end of this chapter:} \emph{I can say a carp, a cicada, a dragonfly, a potato and burdock root.}

Complete the last lesson of chapter 65: i can say a carp, a cicada, a dragonfly, a potato and burdock root.
\end{chapteropening}

\section[koi]{\ja{こい} (koi) — a carp}

\label{lesson:JA-C65-koi}



\begin{quote}
Before the new one: say the Japanese for a sardine, then the Japanese for a flavour.
\end{quote}

\subsection*{You'll want to know: \ja{こい}}
\textbf{\ja{こい}} — \emph{koi} — \textquotedblleft{}a carp\textquotedblright{}.

\begin{grammarlens}[title={What you've built}]
One more word for this chapter.
\end{grammarlens}

\subsection*{Your turn}
\begin{itemize}
\raggedright
  \item \emph{Say it:} \emph{koi}
  \item \emph{Say it:} \emph{koi}, once more
  \item \emph{Say it:} say \emph{koi}
  \item \emph{Recall:} say the Japanese for a sardine, then the Japanese for a flavour, then say \emph{koi} again
\end{itemize}

\begin{tcolorbox}[breakable,colback=teal!4,colframe=teal!35!black,title={Before you move on}]
What does \ja{こい} mean? (\textquotedblleft{}A carp\textquotedblright{}.) Say it once more.
\end{tcolorbox}
\section[semi]{\ja{せみ} (semi) — a cicada}

\label{lesson:JA-C65-semi}



\begin{quote}
Before the new one: say the Japanese for a flavour, then the Japanese for a carp.
\end{quote}

\subsection*{You'll want to know: \ja{せみ}}
\textbf{\ja{せみ}} — \emph{semi} — \textquotedblleft{}a cicada\textquotedblright{}.

\begin{grammarlens}[title={What you've built}]
One more word for this chapter.
\end{grammarlens}

\subsection*{Your turn}
\begin{itemize}
\raggedright
  \item \emph{Say it:} \emph{semi}
  \item \emph{Say it:} \emph{semi}, once more
  \item \emph{Say it:} say \emph{semi}
  \item \emph{Recall:} say the Japanese for a flavour, then the Japanese for a carp, then say \emph{semi} again
\end{itemize}

\begin{tcolorbox}[breakable,colback=teal!4,colframe=teal!35!black,title={Before you move on}]
What does \ja{せみ} mean? (\textquotedblleft{}A cicada\textquotedblright{}.) Say it once more.
\end{tcolorbox}
\section[tonbo]{\ja{とんぼ} (tonbo) — a dragonfly}

\label{lesson:JA-C65-tonbo}



\begin{quote}
Before the new one: say the Japanese for a carp, then the Japanese for a cicada.
\end{quote}

\subsection*{You'll want to know: \ja{とんぼ}}
\textbf{\ja{とんぼ}} — \emph{tonbo} — \textquotedblleft{}a dragonfly\textquotedblright{}.

\begin{grammarlens}[title={What you've built}]
One more word for this chapter.
\end{grammarlens}

\subsection*{Your turn}
\begin{itemize}
\raggedright
  \item \emph{Say it:} \emph{tonbo}
  \item \emph{Say it:} \emph{tonbo}, once more
  \item \emph{Say it:} say \emph{tonbo}
  \item \emph{Recall:} say the Japanese for a carp, then the Japanese for a cicada, then say \emph{tonbo} again
\end{itemize}

\begin{tcolorbox}[breakable,colback=teal!4,colframe=teal!35!black,title={Before you move on}]
What does \ja{とんぼ} mean? (\textquotedblleft{}A dragonfly\textquotedblright{}.) Say it once more.
\end{tcolorbox}
\section[imo]{\ja{いも} (imo) — a potato}

\label{lesson:JA-C65-imo}



\begin{quote}
Before the new one: say the Japanese for a cicada, then the Japanese for a dragonfly.
\end{quote}

\subsection*{You'll want to know: \ja{いも}}
\textbf{\ja{いも}} — \emph{imo} — \textquotedblleft{}a potato\textquotedblright{}.

\begin{grammarlens}[title={What you've built}]
One more word for this chapter.
\end{grammarlens}

\subsection*{Your turn}
\begin{itemize}
\raggedright
  \item \emph{Say it:} \emph{imo}
  \item \emph{Say it:} \emph{imo}, once more
  \item \emph{Say it:} say \emph{imo}
  \item \emph{Recall:} say the Japanese for a cicada, then the Japanese for a dragonfly, then say \emph{imo} again
\end{itemize}

\begin{tcolorbox}[breakable,colback=teal!4,colframe=teal!35!black,title={Before you move on}]
What does \ja{いも} mean? (\textquotedblleft{}A potato\textquotedblright{}.) Say it once more.
\end{tcolorbox}
\section[gobou]{\ja{ごぼう} (gobou) — burdock root}

\label{lesson:JA-C65-gobou}



\begin{quote}
Before the new one: say the Japanese for a dragonfly, then the Japanese for a potato.
\end{quote}

\subsection*{You'll want to know: \ja{ごぼう}}
\textbf{\ja{ごぼう}} — \emph{gobou} — \textquotedblleft{}burdock root\textquotedblright{}.

\begin{grammarlens}[title={What you've built}]
One more word for this chapter.
\end{grammarlens}

\subsection*{Your turn}
\begin{itemize}
\raggedright
  \item \emph{Say it:} \emph{gobou}
  \item \emph{Say it:} \emph{gobou}, once more
  \item \emph{Say it:} say \emph{gobou}
  \item \emph{Recall:} say the Japanese for a dragonfly, then the Japanese for a potato, then say \emph{gobou} again
\end{itemize}

\begin{tcolorbox}[breakable,colback=teal!4,colframe=teal!35!black,title={Before you move on}]
What does \ja{ごぼう} mean? (\textquotedblleft{}Burdock root\textquotedblright{}.) Say it once more.
\end{tcolorbox}
